\documentclass[a4paper,12pt]{article}

% ---------------- Packages ----------------
\usepackage[utf8]{inputenc}
\usepackage[T1]{fontenc}
\usepackage[french]{babel}
\usepackage{graphicx}
\usepackage{geometry}
\usepackage{hyperref}
\usepackage{fancyhdr}
\usepackage{listings}
\usepackage{xcolor}
\usepackage{booktabs}

\geometry{margin=2.5cm}

% Configuration du code Python
\lstset{
    language=Python,
    basicstyle=\ttfamily\footnotesize,
    keywordstyle=\color{blue}\bfseries,
    stringstyle=\color{red},
    commentstyle=\color{green!60!black},
    numbers=left,
    numberstyle=\tiny\color{gray},
    stepnumber=1,
    numbersep=10pt,
    backgroundcolor=\color{gray!10},
    frame=single,
    rulecolor=\color{gray!30},
    tabsize=4,
    captionpos=b,
    breaklines=true
}

% ---------------- En-tête et pied de page ----------------
\pagestyle{fancy}
\fancyhf{}
\fancyhead[L]{Mini-projet Cybersécurité}
\fancyhead[R]{Analyseur de Sécurité}
\fancyfoot[C]{\thepage}

% ---------------- Infos ----------------
\title{\textbf{MiniSec Scanner}}
\author{
Étudiant en Cybersécurité -- Matricule : 23002su \\
Promotion 2024-2025
}
\date{\today}

\begin{document}

% ---------------- Page de garde ----------------
\maketitle
\vspace{1cm}

\begin{center}
\textbf{Option :} Sécurité Informatique \\
\textbf{Matière :} Cybersécurité \\
\textbf{Outil :} Python 3.x + YARA
\end{center}

\newpage

% ---------------- Résumé ----------------
\section*{Résumé}
Ce mini-projet consiste à développer un outil d'analyse de sécurité de fichiers
appelé MiniSec Scanner. L'objectif est de détecter des fichiers potentiellement
suspects en utilisant des techniques d'analyse heuristique combinées avec le
moteur de détection YARA pour une identification avancée de menaces.

\newpage

% ---------------- Introduction ----------------
\section{Introduction}
La sécurité informatique moderne nécessite des outils efficaces pour détecter
les menaces potentielles. MiniSec Scanner est un outil pédagogique conçu pour
analyser des fichiers et identifier les risques de sécurité basé sur plusieurs
critères : extension, nom, taille et contenu.

\section{Objectif du projet}
L'objectif de ce travail est de :
\begin{itemize}
\item Concevoir un analyseur de fichiers multi-critères.
\item Intégrer le moteur YARA pour une détection professionnelle.
\item Générer des rapports détaillés d'analyse.
\item Respecter les bonnes pratiques de développement Python.
\end{itemize}

% ---------------- Architecture ----------------
\section{Architecture du Scanner}
L'architecture étudiée est basée sur une analyse modulaire.
Elle comprend :
\begin{itemize}
\item Une classe principale \texttt{MiniSecScanner}.
\item Des modules d'analyse (extension, nom, taille, contenu).
\item Un moteur YARA intégré avec signatures embarquées.
\item Un générateur de rapports textuels.
\end{itemize}

L'analyse est réalisée sur des dossiers locaux avec option de scan récursif
afin d'observer les différents niveaux de risque.

% ---------------- Implémentation ----------------
\section{Implémentation en Python}
L'implémentation est réalisée en Python 3.x avec une architecture orientée objet.
Les principaux composants analysés sont :
\begin{itemize}
\item Les extensions dangereuses (.exe, .bat, .dll, etc.).
\item Les mots-clés suspects (virus, malware, crack, etc.).
\item Les anomalies de taille (vide, très volumineux).
\item Les signatures YARA pour détection avancée.
\end{itemize}

\begin{lstlisting}[caption={Exemple d'analyse d'extension}]
def analyze_extension(self, extension: str) -> Tuple[int, str]:
    if extension in self.HIGH_RISK_EXTENSIONS:
        return (3, f"Extension dangereuse : {extension}")
    elif extension in self.MEDIUM_RISK_EXTENSIONS:
        return (2, f"Extension suspecte : {extension}")
    return (0, "")
\end{lstlisting}

% ---------------- YARA Integration ----------------
\section{Intégration du Moteur YARA}
YARA est intégré comme dépendance optionnelle pour une détection professionnelle.
Les règles sont embarquées directement dans le code :

\begin{lstlisting}[caption={Règles YARA intégrées}]
YARA_RULES = """
rule EICAR_Test_File {
    meta:
        description = "Standard antivirus test file"
        threat_level = "Medium"
    strings:
        $eicar = "EICAR-STANDARD-ANTIVIRUS-TEST-FILE"
    condition:
        $eicar
}
"""
\end{lstlisting}

% ---------------- Résultats ----------------
\section{Résultats et analyse}
Les résultats obtenus montrent que le scanner détecte efficacement les menaces.
Les fichiers sont classifiés en trois niveaux de risque :
\begin{itemize}
\item \textbf{Low} : Fichiers standards (0-1 point)
\item \textbf{Medium} : Extensions suspectes (2-4 points)
\item \textbf{High} : Menaces claires (5+ points)
\end{itemize}

\begin{table}[h]
\centering
\begin{tabular}{|l|c|c|}
\hline
\textbf{Type de fichier} & \textbf{Risque} & \textbf{Score} \\
\hline
document.txt & Low & 0 \\
program.exe & Medium & 2 \\
crack\_tool.exe & High & 6 \\
eicar\_test.txt & Medium & 3 \\
\hline
\end{tabular}
\caption{Résultats d'analyse sur fichiers de test}
\end{table}

% ---------------- Conclusion ----------------
\section{Conclusion}
Ce mini-projet a permis de comprendre les principes de l'analyse de sécurité
et de développer un outil fonctionnel avec des technologies professionnelles.
Les résultats confirment que l'approche heuristique combinée à YARA offre
une détection efficace tout en restant pédagogique.

% ---------------- Déclaration d'honneur ----------------
\section*{Déclaration d'honneur}
Je déclare que ce travail a été réalisé par mes soins et qu'il est
exempt de tout plagiat.

\end{document}
